\chapter{Mở đầu}
\label{chap:modau}
\section{Bối cảnh}
\label{sec:boi-canh}
\subsection{Giới hạn của phương pháp số cổ điển}
\label{subsec:gioi-han-co-dien}
Việc mô hình hoá và dự báo động lực học của các hệ đa vật lý, đa quy mô vẫn là
một bài toán mở của khoa học hiện đại. Các hệ thực tế thường chịu chi phối đồng
thời của nhiều quá trình vật lý, hoá học và sinh học tương tác trên những thang
không--thời gian rất khác nhau. Hệ Trái Đất là ví dụ điển hình: các quá trình chi
phối nó trải rộng trên $17$ bậc độ lớn về thang không--thời
gian~\cite{karniadakis2021}.
 
Trong nhiều thập kỷ, việc giải số các phương trình đạo hàm riêng đã đóng vai trò
trung tâm trong nghiên cứu và mô phỏng các hệ vật lý. Các phương pháp cổ điển --
sai phân hữu hạn, phần tử hữu hạn, phương pháp phổ và các phương pháp không lưới
-- đã đạt nhiều thành công từ địa vật lý, cơ học chất lưu đến lý sinh và khoa học
vật liệu. Tuy nhiên, với các hệ phi tuyến, đa quy mô và có tầng bậc thang không
đồng nhất, các phương pháp này vẫn vấp phải chi phí tính toán rất lớn và nhiều
nguồn bất định.
 
Ba giới hạn cụ thể đáng được nêu riêng, vì cả ba đều là động lực trực tiếp cho
hướng tiếp cận của báo cáo~\cite{karniadakis2021}.
 
\emph{Thứ nhất, bài toán ngược.} Khi mục tiêu không chỉ là tìm nghiệm của một
phương trình đã biết mà còn phải ước lượng tham số chưa biết, tính chất vật liệu,
hoặc chính cơ chế vật lý còn thiếu, các phương pháp cổ điển đòi hỏi công thức hoá
riêng, thuật toán mới và mã nguồn phức tạp. Chi phí thường ở mức không chấp nhận
được, vì mỗi bước lặp của bài toán ngược lại đòi hỏi giải trọn vẹn một bài toán
thuận.
 
\emph{Thứ hai, điều kiện biên khuyết hoặc nhiễu.} Các phương pháp số cổ điển giả
thiết điều kiện đầu và điều kiện biên được cho đầy đủ và chính xác. Trong nhiều
hệ thực tế, các thông tin này bị thiếu, gián đoạn hoặc nhiễu; khi đó việc áp dụng
trực tiếp các phương pháp truyền thống là bất khả thi.
 
\emph{Thứ ba, lời nguyền số chiều.} Với các bài toán nhiều chiều, mọi phương pháp
dựa trên lưới đều chịu chi phí tăng theo hàm mũ của số chiều. Đây là giới hạn có
bản chất toán học chứ không phải giới hạn của cài đặt, và
Nhận xét~\ref{nx:loi-nguyen} của Chương~\ref{chap:fnn} sẽ định lượng nó bằng cách
đối chiếu bậc xấp xỉ $\mathcal{O}(m^{-s/d})$ của các phương pháp tuyến tính cổ
điển với bậc $\mathcal{O}(m^{-1/2})$ không phụ thuộc số chiều của mạng nơ-ron.
 
\subsection{Dữ liệu quan sát và giới hạn của học máy thuần dữ liệu}
 
Song song với những khó khăn trên, sự phát triển của hệ thống cảm biến, thiết bị
đo đạc và công nghệ viễn thám đã tạo ra khối lượng dữ liệu quan sát chưa từng có.
Dữ liệu đến từ cảm biến mặt đất, thiết bị quan trắc trên biển và trên không, cùng
các hệ thống vệ tinh; chúng có độ phân giải, độ chính xác và mức tin cậy khác
nhau, tạo thành các tập dữ liệu \emph{đa độ trung thực} với phạm vi bao phủ rộng
trên cả không gian lẫn thời gian~\cite{karniadakis2021}.
 
Sự dồi dào ấy mở ra cơ hội cho các phương pháp dựa trên dữ liệu. Học máy có khả
năng khai thác không gian dữ liệu nhiều chiều, phát hiện tương quan phức tạp và
xử lý những bài toán đặt không chỉnh; học sâu nói riêng cung cấp công cụ tự động
trích xuất đặc trưng từ dữ liệu có cấu trúc phức tạp. Các phương pháp này đã được
dùng để phát hiện hiện tượng khí hậu cực đoan, dự báo lượng mưa và phân tích dữ
liệu trong khoa học vật liệu.
 
Tuy nhiên, mô hình học máy \emph{thuần dữ liệu} có một giới hạn cơ bản khi áp dụng
cho hệ vật lý: nó có thể đạt sai số nhỏ trên tập quan sát mà vẫn đưa ra dự đoán
vô lý về mặt vật lý khi ngoại suy, đơn giản vì không có gì ràng buộc nó tuân theo
các định luật bảo toàn hay cấu trúc toán học chi phối hệ. Thêm vào đó, trong nhiều
bài toán khoa học, dữ liệu quan sát thưa, không đồng nhất và chứa nhiễu -- đúng
\emph{chế độ ít dữ liệu} mà Raissi và cộng sự~\cite{raissi2019} nêu là nơi các
công cụ học sâu tiêu chuẩn thiếu tính vững và không có bảo đảm hội tụ nào.
 
\subsection{Nhúng vật lý vào quá trình học}
 
Từ hai nhận định đối lập trên -- phương pháp số cổ điển mạnh về cấu trúc nhưng
yếu khi dữ kiện khuyết, học máy mạnh về dữ liệu nhưng mù về cấu trúc -- hướng
tiếp cận tự nhiên là kết hợp cả hai: đưa tri thức tiên nghiệm về hệ vào chính quá
trình học. Tri thức ấy có thể là định luật vật lý, nguyên lý bảo toàn, quan hệ
toán học hay kiến thức chuyên ngành; việc tích hợp nó thu hẹp không gian nghiệm
và hướng quá trình tối ưu về những nghiệm phù hợp quy luật vật lý. Đây là nội
dung của \emph{học máy khoa học}.
 
Karniadakis và cộng sự~\cite{karniadakis2021} phân loại ba con đường nhúng vật lý,
và sự phân loại này định vị chính xác phạm vi của báo cáo.
 
\emph{Thiên kiến quan sát} (observational bias) đưa vật lý vào qua chính dữ liệu:
huấn luyện trên dữ liệu đã hàm chứa quy luật vật lý cho phép mô hình học được các
hàm và trường phản ánh cấu trúc ấy. Cách này đơn giản về khái niệm nhưng đòi hỏi
khối lượng dữ liệu lớn, mà chi phí thu thập trong khoa học và kỹ thuật thường rất
cao.
 
\emph{Thiên kiến quy nạp} (inductive bias) đưa vật lý vào qua \emph{kiến trúc}:
thiết kế mạng sao cho mọi dự đoán \emph{đương nhiên} thoả một tập ràng buộc. Đây
là cách nguyên tắc nhất vì ràng buộc được thoả một cách chặt chẽ, nhưng nó thường
chỉ áp dụng được cho các nhóm đối xứng tương đối đơn giản đã biết trước, và dẫn
tới cài đặt phức tạp khó mở rộng.
 
\emph{Thiên kiến học} (learning bias) đưa vật lý vào qua \emph{hàm mục tiêu}: chọn
hàm mất mát và ràng buộc sao cho quá trình huấn luyện ưu tiên các nghiệm tuân theo
vật lý. Bằng các ràng buộc phạt mềm, định luật vật lý chỉ được thoả \emph{xấp
xỉ}; đổi lại, đây là nền linh hoạt nhất, cho phép đưa vào một lớp rất rộng các
ràng buộc dạng tích phân, vi phân hay thậm chí phân thứ.
 
Ba con đường này không loại trừ nhau. Báo cáo tập trung vào con đường thứ ba, và
sự lựa chọn ấy có cái giá của nó: Nhận xét~\ref{rem:soft-constraint} của
Chương~\ref{chap:pinn} sẽ chỉ ra rằng phạt mềm khiến điều kiện biên không được
thoả chính xác, còn Mục~\ref{sec:gradient-pathology} phân tích khó khăn tối ưu
phát sinh khi nhiều mục tiêu cạnh tranh bị gộp vào một hàm mục tiêu duy nhất.
Bổ đề~\ref{bd:rang-buoc-cung} trình bày lựa chọn thay thế theo tinh thần thiên
kiến quy nạp, để đối chiếu.
 
\subsection{Mạng nơ-ron thông tin vật lý}
 
Đại diện tiêu biểu của con đường thứ ba là \emph{mạng nơ-ron thông tin vật lý}
(physics-informed neural network, PINN). Ý tưởng cốt lõi rất gọn: biểu diễn nghiệm
chưa biết bằng một mạng nơ-ron, rồi đưa chính phương trình chi phối vào hàm mất
mát dưới dạng \emph{phần dư}. Với bài toán tham số hoá dạng

\begin{equation*}
    u_t + \Lop\left[u; \bzeta\right] = 0,
\end{equation*}

theo cách viết của Raissi và cộng sự~\cite{raissi2019}, mạng $u_{\btheta}$ được
huấn luyện để làm nhỏ đồng thời phần dư của phương trình tại một tập điểm trong
miền và sai lệch so với điều kiện đầu, điều kiện biên. Điểm đáng chú ý là các
điểm trong miền \emph{không cần nhãn}: thông tin giám sát tại đó do chính phương
trình cung cấp. Chương~\ref{chap:pinn} xây dựng đầy đủ khung này.
 
Ý tưởng dùng mạng nơ-ron làm hàm thử cho phương trình vi phân không mới; nó đã
xuất hiện từ những năm 1990 với các công trình của Dissanayake và
Phan-Thien~\cite{dissanayake1994} và của Lagaris cùng cộng sự~\cite{lagaris1998}.
Điều làm nó trở nên khả thi ở quy mô hiện nay là hai tiến bộ công cụ: vi phân tự
động, cho phép tính đạo hàm theo biến đầu vào chính xác tới độ chính xác máy mà
không cần lưới; và các khung tính toán khả vi cùng phần cứng hiện đại.
Chương~\ref{chap:autodiff} dành trọn cho hai công cụ này, và
Mệnh đề~\ref{md:fd-rao-can} sẽ định lượng vì sao sai phân hữu hạn \emph{không}
thay thế được vi phân tự động trong bối cảnh này.
 
PINNs đặc biệt phù hợp với các bài toán mà vật lý và dữ liệu chỉ được biết một
phần. Ở một cực, phương trình và tham số đã xác định đầy đủ, dữ liệu chỉ cung cấp
điều kiện đầu hoặc điều kiện biên -- đó là \emph{bài toán thuận}. Ở cực kia, có
nhiều dữ liệu quan sát nhưng chưa rõ phương trình chi phối. Phần lớn bài toán thực
tế nằm giữa: một phần tri thức vật lý đã biết, còn một số tham số hoặc quan hệ cấu
thành cần được suy luận từ dữ liệu -- đó là \emph{bài toán ngược}. Cùng một khung
hàm mất mát phục vụ được cả hai, vì tham số vật lý được xử lý hoàn toàn như tham
số của mạng; Mục~\ref{subsec:thuan-nguoc} trình bày điều này.
 
Cuối cùng, cần nói ngay rằng PINNs không phải một phương pháp không có nhược điểm,
và chính các nhược điểm ấy chiếm phần lớn dung lượng của báo cáo. Karniadakis và
cộng sự~\cite{karniadakis2021} nêu thẳng rằng mạng kết nối đầy đủ gặp khó khăn khi
học các hàm tần số cao -- hiện tượng gọi là \emph{thiên kiến phổ} -- và rằng các
đặc trưng tần số cao trong nghiệm sinh ra gradient dốc, khiến PINNs thường không
phạt được chính xác phần dư. Chương~\ref{chap:pinn} truy nguyên hiện tượng này tới
ba nguyên nhân độc lập ở ba tầng khác nhau (Bảng~\ref{tab:ba-nguyen-nhan}), và
Chương~\ref{chap:thucnghiem} kiểm chứng chúng bằng số liệu.
 
%==============================================================================
\section{Mục tiêu nghiên cứu}
\label{sec:muc-tieu}
%==============================================================================
 
\subsection{Mục tiêu tổng quát}
 
Báo cáo hướng tới việc trình bày một cách hệ thống nền tảng toán học, thuật toán
và thực nghiệm của PINNs khi được dùng làm công cụ giải phương trình vi phân;
đồng thời cài đặt và đánh giá định lượng phương pháp này trên các bài toán chuẩn
mực, nhằm xác lập một đánh giá cân bằng -- dựa trên số liệu và chứng minh chứ
không dựa trên kỳ vọng -- về phạm vi mà PINNs thực sự mang lại lợi thế so với giải
tích số cổ điển.
 
Hai chữ ``cân bằng'' cần được hiểu theo nghĩa mạnh. Những kết quả không đạt kỳ
vọng cũng được trình bày đầy đủ, và trong một trường hợp, một kết quả tiêu cực
đã dẫn tới việc phải siết lại cách phát biểu của chính chương lý thuyết
(Nhận xét~\ref{nx:hieu-chinh}).
 
\subsection{Mục tiêu cụ thể}
\label{subsec:mt-cu-the}
 
Từ mục tiêu tổng quát, báo cáo đặt ra năm nhiệm vụ.
 
\begin{enumerate}[label=\textbf{MT\arabic*.},leftmargin=*,itemsep=5pt]
\item \textbf{Hệ thống hoá cơ sở toán học của mạng nơ-ron truyền thẳng.} Trình bày
kiến trúc, tính chất xấp xỉ phổ quát và vai trò của hàm kích hoạt, làm nền tảng
cho việc dùng mạng nơ-ron như một không gian hàm thử thay cho không gian phần tử
hữu hạn hay không gian phổ. Trọng tâm đặt vào một câu hỏi mà tài liệu học sâu tổng
quát bỏ qua: \emph{đạo hàm bậc cao của mạng theo biến đầu vào có dạng tường minh
nào}, và điều kiện nào trên hàm kích hoạt là \emph{cần} để chúng không đồng nhất
triệt tiêu (Chương~\ref{chap:fnn}).
 
\item \textbf{Làm rõ cơ chế vi phân tự động và thuật toán tối ưu.} Phân tích cách
vi phân tự động cho phép tính chính xác tới độ chính xác máy các đạo hàm riêng bậc
cao xuất hiện trong toán tử $\Lop$, và so sánh các thuật toán tối ưu dùng phổ biến
trong huấn luyện PINNs, đặc biệt là cặp Adam và L-BFGS
(Chương~\ref{chap:autodiff}).
 
\item \textbf{Xây dựng phương pháp PINNs một cách chặt chẽ.} Thiết lập đầy đủ công
thức của mô hình thời gian liên tục, từ định nghĩa phần dư \eqref{eq:residual} tới
hàm mục tiêu tổng hợp \eqref{eq:total-loss}; làm rõ điều kiện để việc cực tiểu hoá
phần dư là hợp lý; và phân tích khó khăn của phương pháp ở cả tầng tối ưu lẫn tầng
xấp xỉ (Chương~\ref{chap:pinn}).
 
\item \textbf{Cài đặt và thực nghiệm trên PyTorch.} Xây dựng một cài đặt PINNs
hoàn chỉnh, minh bạch về mã nguồn, và áp dụng cho các bài toán cụ thể; báo cáo sai
số $L^2$ tương đối, sai số $L^{\infty}$, đường cong hội tụ và chi phí tính toán
thực đo (Chương~\ref{chap:thucnghiem}).
 
\item \textbf{Khảo sát độ nhạy và phân tích hạn chế.} Đánh giá ảnh hưởng của hàm
kích hoạt, trọng số trong hàm mất mát và lựa chọn thuật toán tối ưu tới độ chính
xác; đối chiếu từng dự báo lý thuyết với số liệu đo được trong
Bảng~\ref{tab:doi-chieu}, và nêu thẳng những hạn chế của chính thiết kế thực
nghiệm ở Mục~\ref{sec:hanche}.
\end{enumerate}
 
\subsection{Phạm vi và giới hạn}
\label{subsec:pham-vi}
 
Để bảo đảm chiều sâu, báo cáo giới hạn phạm vi như sau.
 
Về \emph{lớp bài toán}, trọng tâm đặt vào bài toán thuận; bài toán ngược chỉ được
trình bày ở mức tổng quan (Mục~\ref{subsec:thuan-nguoc}) nhằm cho thấy tính thống
nhất của khung phương pháp.
 
Về \emph{cơ chế nhúng vật lý}, báo cáo tập trung vào thiên kiến học, tức cưỡng
bức bằng phạt mềm; thiên kiến quy nạp chỉ được nhắc tới khi cần đối chiếu, qua
dạng nghiệm thử ràng buộc cứng ở Bổ đề~\ref{bd:rang-buoc-cung}.
 
Về \emph{mô hình thời gian}, báo cáo chỉ xét mô hình thời gian liên tục của
Raissi và cộng sự. Mô hình thời gian rời rạc dựa trên các lược đồ Runge--Kutta ẩn
nằm ngoài phạm vi.
 
Về \emph{kiến trúc}, chỉ xét mạng truyền thẳng kết nối đầy đủ. Các họ kiến trúc
khác -- mạng tích chập, mạng đồ thị, hay các toán tử nơ-ron như DeepONet và Fourier
neural operator -- không được đề cập, ngoại trừ kiến trúc có cổng ở
Mục~\ref{subsec:gated}, vốn được đưa vào vì nó là một phần của phân tích bệnh lý
gradient chứ không phải như một kiến trúc độc lập.
 
Về \emph{công cụ}, toàn bộ phần thực nghiệm được cài đặt bằng PyTorch, không dùng
thư viện PINNs đóng gói sẵn, nhằm giữ cho mọi thành phần của thuật toán đều minh
bạch và kiểm chứng được.
 
Cuối cùng, về \emph{tham vọng lý thuyết}: báo cáo không đặt mục tiêu chứng minh
định lý mới về hội tụ hay ước lượng sai số cho PINNs, và các kết quả lớn của lĩnh
vực chỉ được trích dẫn cùng diễn giải. Tuy vậy, báo cáo \emph{có} cung cấp chứng
minh đầy đủ cho một số phát biểu mà tài liệu thường chỉ nêu định tính -- chẳng hạn
sự sụp đổ gradient của kiến trúc ReLU (Hệ quả~\ref{hq:pinn-hong}), chặn ổn định
nối phần dư với sai số (Mệnh đề~\ref{md:chan-on-dinh}), sàn nhiễu của bước học
hằng số (Mệnh đề~\ref{md:san-nhieu}), và chế độ phản hồi dương của quy tắc cân
bằng trọng số (Bổ đề~\ref{bd:phan-hoi-duong}). Các chứng minh này sơ cấp và không
phải đóng góp lý thuyết mới của lĩnh vực; giá trị của chúng nằm ở chỗ biến những
khẳng định mơ hồ thành những phát biểu có giả thiết rõ ràng và kiểm chứng được.
 
\subsection{Đóng góp dự kiến}
 
Trong phạm vi đã nêu, báo cáo có ba đóng góp.
 
\emph{Thứ nhất}, một trình bày mạch lạc bằng tiếng Việt về PINNs, đi từ cơ sở toán
học của mạng nơ-ron truyền thẳng tới công thức đầy đủ của mô hình thời gian liên
tục, với notation nhất quán xuyên suốt và mọi khẳng định đều truy được về một mệnh
đề cụ thể. Đây là nội dung hiện còn tương đối thiếu trong tài liệu học thuật tiếng
Việt.
 
\emph{Thứ hai}, một cài đặt PyTorch độc lập kèm kết quả thực nghiệm định lượng,
tái lập được, có thể dùng làm điểm khởi đầu cho các nghiên cứu tiếp theo.
 
\emph{Thứ ba}, một phân tích đối chiếu có hệ thống giữa lý thuyết và số liệu
(Bảng~\ref{tab:doi-chieu}), qua đó chỉ ra cụ thể những tình huống PINNs phát huy
hiệu quả và những tình huống phương pháp gặp khó khăn. Phần đối chiếu này cho một
kết luận phương pháp luận đáng chú ý: \emph{không} một mệnh đề đã chứng minh nào
bị số liệu bác bỏ --- mọi chỗ phải siết lại đều đến từ suy luận tiệm cận, từ
việc ghép hai kết quả rời rạc, từ một ước lượng chưa kiểm chứng, hoặc từ việc
khái quát hoá một lần chạy duy nhất.
 
%==============================================================================
\section{Bố cục báo cáo}
\label{sec:bo-cuc}
%==============================================================================
 
Bốn chương còn lại được tổ chức theo một mạch nối tiếp: mỗi chương chỉ dùng kết
quả của các chương trước, và mỗi lựa chọn thiết kế ở chương sau đều truy được về
một mệnh đề ở chương trước.
 
\textbf{Chương~\ref{chap:fnn}} thiết lập cơ sở toán học của mạng truyền thẳng, với
trọng tâm khác hẳn tài liệu học sâu tổng quát: đạo hàm của mạng theo \emph{biến
đầu vào}. Hai hệ thức truy hồi cho Jacobi và Hessian
(Mệnh đề~\ref{md:jacobian}, \ref{md:hessian}) là công cụ trung tâm cho mọi tính
toán phần dư về sau. Chương kết thúc bằng một cấu hình kiến trúc được biện minh
từng thành phần -- hàm kích hoạt $\tanh$, khởi tạo Xavier, độ chệch bằng không,
lớp ra tuyến tính -- dùng cố định cho toàn bộ phần thực nghiệm.
 
\textbf{Chương~\ref{chap:autodiff}} trình bày hai công cụ tính toán: vi phân tự
động và các thuật toán tối ưu. Kết quả then chốt là
Mệnh đề~\ref{md:ad-tai-tao}, chứng minh rằng vi phân tự động tái tạo \emph{đúng}
hai hệ thức truy hồi của Chương~\ref{chap:fnn} mà không cần lập trình tường minh
đạo hàm bậc hai của hàm kích hoạt. Chiến lược huấn luyện hai giai đoạn
Adam~$\to$~L-BFGS được biện minh bằng ba lập luận cụ thể chứ không bằng kinh
nghiệm thực hành.
 
\textbf{Chương~\ref{chap:pinn}} xây dựng phương pháp PINNs theo ba tầng: cấu trúc
hàm mất mát, khó khăn ở tầng tối ưu, và giới hạn ở tầng xấp xỉ. Đóng góp tổng hợp
của chương là Bảng~\ref{tab:ba-nguyen-nhan}: ba nguyên nhân độc lập, nằm ở ba tầng
khác nhau, cùng làm PINNs suy giảm trên đúng lớp bài toán có thang độ dài nhỏ.
 
\textbf{Chương~\ref{chap:thucnghiem}} cài đặt toàn bộ khung ấy trên PyTorch và
kiểm chứng bằng mười thí nghiệm, mỗi thí nghiệm trả lời một câu hỏi cụ thể với tiêu
chí thành công xác định trước. Kết quả được tổng hợp ở Bảng~\ref{tab:doi-chieu},
và những hạn chế của chính thiết kế thực nghiệm được nêu thẳng ở
Mục~\ref{sec:hanche}.